% !TEX program = xelatex
\documentclass[UTF8,11pt]{ctexart}

\usepackage[a4paper,margin=2.2cm]{geometry}
\usepackage{hyperref}
\usepackage{bookmark}
\usepackage{enumitem}
\usepackage{tabularx}
\usepackage{array}
\usepackage{booktabs}
\usepackage{xcolor}
\usepackage{xurl}
\usepackage{longtable}

\hypersetup{
  colorlinks=true,
  linkcolor=blue,
  urlcolor=blue,
  citecolor=blue
}

\newcolumntype{Y}{>{\raggedright\arraybackslash}X}
\setlength{\emergencystretch}{2em}

\title{OpenAI Codex 最新主线与当前定制版差距分析报告}
\author{jqwang / Codex}
\date{2026年3月15日}

\begin{document}
\maketitle

\noindent\textbf{状态说明（2026-03-24 更新）}：本文是当时的差距分析快照。其中涉及的 macOS Endpoint Security、本地 Smart Access、\texttt{/freeze} 等内容，已不属于当前 upstream merge 目标。当前方向是移除本地 \texttt{endpoint-sec} / \texttt{request\_security\_override} / \texttt{freeze} 方案，保留账号池等定制能力，并继续向官方 \texttt{approval\_policy} + \texttt{approvals\_reviewer} / guardian 语义收敛。

\begin{abstract}
本文基于 \href{https://github.com/openai/codex}{OpenAI Codex 官方仓库}当前最新主线，以及本地当前定制版 \texttt{main} 分支的实际代码状态，对两者的功能差距、架构漂移和合并风险进行一次决策导向的梳理。

结论不是“官方更先进，所以直接升级即可”，而是：\textbf{两边已经分叉成两条产品线}。官方在 app-server v2、插件/应用生态、审批治理、实时链路、工具协议标准化等通用平台能力上明显领先；当前定制版则在 macOS 本机能力、定制安全放行、工作流型子代理、多模型路由、Entire 集成和定制 TUI 体验上更强。未来若要缩小差距，应按功能族逐步迁移，而不应直接把整条分支硬合并。
\end{abstract}

\tableofcontents
\newpage

\section{执行摘要}

\subsection{一句话结论}
\textbf{当前定制版并不是“单纯落后几个版本”，而是与官方最新主线发生了大规模双向分叉。}

\subsection{核心判断}
\begin{itemize}[nosep]
  \item 从主线演进看：官方已经明显向“标准化平台 + app-server 协议 + 审批治理 + 插件/应用生态”继续推进。
  \item 从差异结构看：当前定制版保留了多项官方没有的高价值特性，尤其是面向 macOS、本地安全控制、工作流型协作和多模型运营侧能力。
  \item 从集成成本看：直接整分支合并的风险很高，冲突会集中在四大区域：\texttt{core}、\texttt{tui}、\texttt{app-server}、\texttt{app-server-protocol}。
  \item 从策略上看：更合理的路径是“以最新上游为底座，按功能族回迁定制能力”，而不是继续在当前大分叉分支上滚动同步。
\end{itemize}

\section{比较基线与方法}

\subsection{比较对象}
\begin{itemize}[nosep]
  \item 官方仓库：\url{https://github.com/openai/codex}
  \item 本地定制版：当前工作区 \texttt{main}
  \item 官方最新主线基线：\texttt{FETCH\_HEAD}，对应上游 \texttt{main} 最新提交
\end{itemize}

\subsection{为什么本报告以 \texttt{main} 为主，不以 release tag 为主}
\begin{itemize}[nosep]
  \item 截至 \textbf{2026年3月15日}，GitHub Releases 页面顶部显示最新预发布版本。
  \item 该版本号为 \texttt{rust-v0.115.0-alpha.24}。
  \item 但本次要分析的是“官方现在已经走到哪里”，因此以官方 \texttt{main} 主线更能反映真实功能差距。
  \item 官方当前主线最新提交为 \href{https://github.com/openai/codex/commit/49edf311ac3ae84659b0ec5eacd5e471c881eee8}{\texttt{49edf311}}，提交时间为 \textbf{2026-03-14 22:24:13 -0700}。
\end{itemize}

\subsection{量化基线}
\begin{center}
\small
\begin{tabularx}{\textwidth}{>{\ttfamily}l Y Y}
\toprule
对象 & 标识 & 说明 \\
\midrule
官方主线最新提交 & 49edf311 & 2026-03-14 22:24:13 -0700，提交标题为 \texttt{[apps] Add tool call meta. (\#14647)} \\
当前定制版主线 & 112cd1b4 & 2026-03-15 18:26:42 +0800，提交标题为 \texttt{Merge branch 'feature/upstream-sync'} \\
共同祖先 & 79d6f80e & 2026-02-26 14:11:10 +0000，说明两边自 2 月底之后已发生大规模双向演进 \\
\bottomrule
\end{tabularx}
\end{center}

\begin{itemize}[nosep]
  \item 自共同祖先以来，官方主线独有提交约 \textbf{496} 个。
  \item 自共同祖先以来，当前定制版独有提交约 \textbf{175} 个。
  \item 直接内容差异约为：\textbf{1587 个文件变更，+215149 / -129182 行}。
  \item 这些数字只能说明漂移规模，\textbf{不能直接等价为功能差距}，因为当前定制版已经手工吸收、改写或重实现过部分上游能力。
\end{itemize}

\subsection{本报告的方法}
\begin{enumerate}[nosep]
  \item 先用 Git 基线确认真正的最新上游主线；
  \item 再按功能域比较，而不是只看 commit 数；
  \item 对“官方已有、当前定制版尚未吸收”的能力做重点梳理；
  \item 对“当前定制版独有、官方无直接替代”的能力单独保留；
  \item 最后给出迁移优先级，而不是给出误导性的“建议直接同步”。
\end{enumerate}

\section{官方主线相比当前定制版的主要新增能力}

\subsection{App Server v2 的文件系统 RPC 更完整}
官方主线已经补齐 app-server v2 的文件系统操作面，包括 \texttt{fs/readFile}、\texttt{fs/writeFile}、\texttt{fs/createDirectory}、\texttt{fs/readDirectory}、\texttt{fs/copy}、\texttt{fs/remove} 等。这意味着官方正在把 app-server 明确推进为一个更标准的“远程编排接口层”，而不是仅服务于本地 TUI。

当前定制版里缺少对应的本地实现文件，例如本地并不存在 \path{codex-rs/app-server/src/fs_api.rs}。这类差距不是小修小补能补平的，而是接口面和产品定位已经发生了变化。

\subsection{Plugin / App 生态继续扩展}
官方主线已经有更完整的 \texttt{plugin/read} 能力和对应 schema，说明插件/应用发现、元数据读取、能力暴露这条链路正在继续扩充。官方最新主线里可以直接看到 \texttt{plugin/read} 出现在 app-server protocol 中，而当前定制版尚未接上这条协议面。

如果未来要做更完整的 skills / plugin / app 混合生态，官方方向显然走得更远。

\subsection{动态工具可见性继续标准化}
官方已经将动态工具“是否暴露给上下文”的能力做进协议，字段名是 \texttt{exposeToContext}。这类能力的价值在于：工具不再只是“能不能调用”，而是可以更精细地控制“模型当前能不能感知到这个工具”。

当前定制版没有这套协议字段和行为链路，因此在工具编排的可治理性上略显粗糙。

\subsection{TUI 与 app-server 的一体化继续增强}
官方主线中已经存在 \texttt{codex-rs/app-server/src/in\_process.rs}，表明它进一步向“在 TUI / app / runtime 里内嵌 app-server”推进。配合 richer initialize platform fields，这说明上游在为多前端、多宿主、多进程形态做一致化。

当前定制版没有这条集成线，因此如果目标是和官方 app 生态对接，后续会越来越吃亏。

\subsection{Smart Approvals / GuardianSubagent 审批流}
这是本次最值得关注的新能力之一。官方主线已经将审批流从“纯人工弹框确认”推进到“由专门审查子代理做风险判断再决定是否放行”。

其核心逻辑不是放宽权限，而是把 \texttt{on-request} 审批交给一个只读、不可再次提权的 guardian 子代理进行风险评估。它会：
\begin{itemize}[nosep]
  \item 重建与该审批有关的上下文；
  \item 审查精确的计划动作 JSON；
  \item 输出 \texttt{low / medium / high} 风险和 \texttt{0-100} 风险分；
  \item 在风险分低于阈值时自动批准，否则拒绝或继续要求人工确认。
\end{itemize}

这套设计的产品意义很强：\textbf{它不是提权，而是减少无意义的人类确认次数，同时把风险门槛做成可审计的结构化过程。}

\subsection{Realtime v2 / Transcription / Handoff 继续推进}
官方主线在 Realtime v2 parser、transcription 和 handoff 相关链路上仍在快速推进。当前定制版并未对齐这一轮更新，因此如果未来目标包括更标准的实时语音、实时会话编排或语音接管能力，当前定制版会明显落后。

\subsection{Provider 配置治理更规范}
官方主线对 provider / config 的治理比当前定制版更标准化，例如更明确的 \texttt{openai\_base\_url} 配置边界、内建 provider 不应被随意覆盖等。当前定制版在 provider 路线上更偏“运营型定制”，灵活但不够统一。

这意味着：当前定制版更适合内部定制与多通道运营，但官方更适合做长期可维护的平台能力。

\subsection{多代理协议正在收敛为更稳定的上游契约}
官方主线在 multi-agent prompt、handoff、事件元数据、角色约束上持续迭代。当前定制版并不是没有多代理，而是有一套自己的扩展语义，因此这里不是“有无”的差距，而是\textbf{契约已经分叉}。

从长期维护看，越晚对齐，后续 app-server v2、TUI、多前端和自动化工具链的兼容成本越高。

\section{当前定制版相比官方主线的独有优势}

\subsection{历史：macOS Endpoint Security + 人工批准式安全放行}
这是当时定制版最有代表性的差异化能力之一，围绕 macOS Endpoint Security 的本机安全拦截与放行工作流展开。其意义不在于“多一个工具”，而在于把系统级安全事件与代理执行流程联动起来。

这类能力属于强平台绑定、本机控制链路型能力，官方主线目前没有等价替代。

\subsection{历史：\texttt{/freeze} 时间冻结 / 快照调试工作流}
这是当时定制版引入的一条实验性调试路径：通过 \texttt{/freeze} 将问题现场冻结成可分析快照，用于回放和排障。这类能力更接近“复杂代理工作流的法医式调试工具”，官方主线没有直接同类物。

\subsection{子代理隔离 worktree 的 lease / resume 机制}
当前定制版对子代理和线程切分的处理更偏工程执行流：不仅能 spawn 子代理，还可以把它们绑定到独立 worktree，并对 worktree 与 thread 的关系做 lease / resume 持久化。

这对复杂多分支、多任务并行开发非常有价值，而官方主线当前更偏通用协作协议，不强调这类本地 Git 工作流深度绑定。

\subsection{\texttt{model\_sub} / \texttt{utility\_model} 多模型路由}
当前定制版有明显强化过的多模型调度能力，包括：
\begin{itemize}[nosep]
  \item \texttt{model\_sub} / \texttt{utility\_model}；
  \item provider account pool；
  \item auth key fallback；
  \item 多轮 failover；
  \item 对 Antigravity、Gemini、Grok、Gemma 等 provider / model 的扩展支持。
\end{itemize}

这使当前定制版在“多模型运营和资源池调度”上比官方更灵活，但这也是它与上游代码结构冲突最重的部分之一。

\subsection{Entire 集成、摘要生成与 Context Packet}
当前定制版将 Entire 相关工作流、AI 摘要生成、上下文封包等能力做成了一套较完整的增强链路。这种能力更偏内部工作流优化，强调的是信息汇总、上下文压缩和后续协作效率。

官方主线目前没有等价的 Entire 集成，因此这里属于当前定制版的真实差异化资产。

\subsection{MemoryLink 贯通到 app-server v2 / hooks}
当前定制版已经把 \texttt{MemoryLink} 传播到 app-server v2、hooks 和事件层。需要注意的是，这里不能简单说“官方没有 memory，定制版有”。更准确的说法是：\textbf{两边都在做 memory，但 wire contract 和集成路径不一样。}

因此后续若要吸收上游 app-server v2 新能力，memory 相关结构大概率需要重新对齐。

\subsection{TUI Git Graph / Session Bar / Ralph Loop}
当前定制版的 TUI 有明显更强的“工作台化”倾向，例如：
\begin{itemize}[nosep]
  \item Git Graph；
  \item Session Bar；
  \item Ralph Loop；
  \item 与历史会话、子代理、仓库状态关联更紧的交互方式。
\end{itemize}

官方主线在通用交互和平台收敛上更强，但当前定制版在“重度工程使用场景”下的体验更有特色。

\section{差距判断：不是简单的新旧关系，而是能力重心不同}

\subsection{官方主线更强的方向}
\begin{itemize}[nosep]
  \item 更标准的 app-server v2 协议面；
  \item 更完整的 plugin / app 生态基础；
  \item 更细的动态工具治理；
  \item 更成熟的审批治理链路；
  \item 更持续演进的 realtime / transcription 能力；
  \item 更适合多前端、多宿主、多客户端的长期平台化。
\end{itemize}

\subsection{当前定制版更强的方向}
\begin{itemize}[nosep]
  \item 更强的 macOS 本机安全与执行控制；
  \item 更强的工程执行工作流；
  \item 更灵活的多模型与账号池调度；
  \item 更深的内部工具与上下文增强；
  \item 更偏“高强度开发者工作台”的 TUI 使用体验。
\end{itemize}

\subsection{真正的风险点}
真正的风险不在于“官方有没有比我们多一些 commit”，而在于：
\begin{itemize}[nosep]
  \item 上游平台能力仍在快速演进；
  \item 当前定制版在若干核心区域深度改过；
  \item 这导致后续每一次同步都更像一次重新集成，而不是普通 rebase。
\end{itemize}

\section{合并风险与冲突热点}

\subsection{高风险冲突区域}
\begin{longtable}{p{0.22\textwidth} p{0.72\textwidth}}
\toprule
区域 & 风险说明 \\
\midrule
\endhead
\path{codex-rs/core/src} & 同时承载模型调度、审批、memory、子代理、sandbox、工具路由等核心行为，是冲突密度最高的区域。 \\
\path{codex-rs/tui/src} & 当前定制版对 TUI 的产品化增强较多，而官方主线也在持续重构交互与 app-server 对接，容易出现行为回归。 \\
\path{codex-rs/app-server/src} & 官方在这里明显前进更多，尤其是文件系统 RPC、plugin/app 链路、进程内启动方式。 \\
\path{codex-rs/app-server-protocol/src} & wire contract 漂移非常明显，任何粗暴合并都可能造成 schema、TS 生成物和客户端行为不一致。 \\
schema 生成物 & 即使业务逻辑合并成功，schema 生成物也容易大面积漂移，放大 review 难度。 \\
\bottomrule
\end{longtable}

\subsection{为什么不建议直接整分支合并}
\begin{itemize}[nosep]
  \item 历史上已经发生过多轮大规模上游同步，旧冲突解决逻辑会继续污染新的同步；
  \item 当前定制版并不只是局部 patch，而是形成了自己的产品逻辑；
  \item 直接整分支 merge 会把“有价值的差异化能力”和“过时的兼容层”混在一起带回主线，难以治理。
\end{itemize}

\section{建议的迁移优先级}

\subsection{优先吸收的官方能力}
\begin{enumerate}[nosep]
  \item app-server v2 文件系统 RPC；
  \item \texttt{plugin/read} 与更完整的 plugin / app 元数据读取；
  \item \texttt{exposeToContext} 这类动态工具可见性治理；
  \item Smart Approvals / GuardianSubagent 审批流；
  \item provider 配置治理和更清晰的 config 边界。
\end{enumerate}

\subsection{应优先保留的当前定制能力}
\begin{enumerate}[nosep]
  \item macOS Endpoint Security + 安全放行；
  \item \texttt{/freeze} 快照调试；
  \item 子代理独立 worktree 与 lease / resume；
  \item \texttt{model\_sub} / \texttt{utility\_model} / account pool / failover；
  \item Entire 集成、上下文封包、定制摘要链路；
  \item Git Graph / Session Bar / Ralph Loop。
\end{enumerate}

\subsection{建议的工程路线}
\begin{enumerate}[nosep]
  \item 以官方最新 \texttt{main} 为新底座建立干净分支；
  \item 按功能族逐块回迁当前定制能力，而不是整分支回放；
  \item 先处理协议和 app-server，再处理多模型和 memory，再处理 sandbox / TUI 深度定制；
  \item 每一族能力都单独验证，不把“看起来能编译”误当作完成。
\end{enumerate}

\section{最终结论}
\textbf{当前定制版和官方最新主线之间，已经不是“一边先进一边落后”的线性关系，而是“官方平台化继续推进，而当前定制版沿着本地工程工作流深度定制”的双轨结构。}

如果目标是：
\begin{itemize}[nosep]
  \item \textbf{对齐官方生态}，就要优先吸收 app-server v2、plugin/app、审批治理、动态工具治理；
  \item \textbf{保住当前产品特色}，就必须把 macOS 安全、工作流型子代理、多模型路由、Entire 集成和定制 TUI 作为明确资产保留；
  \item \textbf{控制长期维护成本}，就不应再把未来同步当成一次次大 merge，而应改为“上游底座 + 定制能力分层回迁”的路线。
\end{itemize}

换言之，\textbf{差距确实已经很大，但并不是单方向的落后；真正需要的是一份 keep / adopt / drop 迁移矩阵，而不是一次粗暴升级。}

\section*{附录：代表性参考链接}
\small
\begin{itemize}[nosep]
  \item 官方仓库：\url{https://github.com/openai/codex}
  \item 官方最新主线提交：\url{https://github.com/openai/codex/commit/49edf311ac3ae84659b0ec5eacd5e471c881eee8}
  \item 共同祖先提交：\url{https://github.com/openai/codex/commit/79d6f80e41806d61b8a5ce7dbaff4bb1d6e38a91}
  \item Smart Approvals：\url{https://github.com/openai/codex/commit/bc24017d64829d0b97b8bc6ed529a389e1e8bc1b}
  \item app-server FS API：\url{https://github.com/openai/codex/commit/f8f82bfc2b558229cc4f7ef6245c474ee8b389c7}
  \item 动态工具可见性：\url{https://github.com/openai/codex/commit/70eddad6b075f26f0f93c66f7ec9a4e49cdadc93}
  \item TUI 内嵌 app-server：\url{https://github.com/openai/codex/commit/9dba7337f21dbc720bd5af70c1628d7c3217f47b}
  \item Provider 配置治理：\url{https://github.com/openai/codex/commit/4b9d5c8c1bdb6d9cfd43570e0b8e88c88b54d823}
  \item 多代理契约收敛：\url{https://github.com/openai/codex/commit/36dfb844277e79793766f96305c9633f90bc043e}
  \item Realtime v2 parser：\url{https://github.com/openai/codex/commit/3e8f47169e523d2004fe4491bd2e29e78f9c6720}
\end{itemize}

\end{document}
